\section{Periodic complexes and Herbrand quotients}
\label{section-periodic-complexes}

\noindent
Of course there is a very general notion of periodic complexes.
We can require periodicity of the maps, or periodicity of the objects.
We will add these here as needed. For the moment we only need
the following cases.

\begin{definition}
\label{definition-periodic-complex}
Let $R$ be a ring.
\begin{enumerate}
\item A {\it $2$-periodic complex} over $R$ is given
by a quadruple $(M, N, \varphi, \psi)$ consisting of
$R$-modules $M$, $N$ and $R$-module maps $\varphi : M \to N$,
$\psi : N \to M$ such that
$$
\xymatrix{
\ldots \ar[r] &
M \ar[r]^\varphi &
N \ar[r]^\psi &
M \ar[r]^\varphi &
N \ar[r] & \ldots
}
$$
is a complex. In this setting we define the {\it cohomology modules}
of the complex to be the $R$-modules
$$
H^0(M, N, \varphi, \psi) = \Ker(\varphi)/\Im(\psi)
\quad\text{and}\quad
H^1(M, N, \varphi, \psi) = \Ker(\psi)/\Im(\varphi).
$$
We say the $2$-periodic complex is {\it exact} if the cohomology
groups are zero.
\item A {\it $(2, 1)$-periodic complex} over $R$ is given
by a triple $(M, \varphi, \psi)$ consisting of an $R$-module $M$ and
$R$-module maps $\varphi : M \to M$, $\psi : M \to M$
such that
$$
\xymatrix{
\ldots \ar[r] &
M \ar[r]^\varphi &
M \ar[r]^\psi &
M \ar[r]^\varphi &
M \ar[r] & \ldots
}
$$
is a complex. Since this is a special case of a $2$-periodic complex
we have its {\it cohomology modules} $H^0(M, \varphi, \psi)$,
$H^1(M, \varphi, \psi)$ and a notion of exactness.
\end{enumerate}
\end{definition}

\noindent
In the following we will use any result proved for $2$-periodic
complexes without further mention for $(2, 1)$-periodic complexes.
It is clear that the collection of $2$-periodic complexes forms a
category with morphisms
$(f, g) : (M, N, \varphi, \psi) \to (M', N', \varphi', \psi')$
pairs of morphisms $f : M \to M'$ and $g : N \to N'$ such
that $\varphi' \circ f = g \circ \varphi$ and $\psi' \circ g = f \circ \psi$.
We obtain an abelian category, with kernels and cokernels as in
Homology, Lemma \ref{homology-lemma-cat-chain-abelian}.

\begin{definition}
\label{definition-periodic-length}
Let $(M, N, \varphi, \psi)$ be a $2$-periodic complex
over a ring $R$ whose cohomology modules have finite length.
In this case we define the {\it multiplicity} of $(M, N, \varphi, \psi)$
to be the integer
$$
e_R(M, N, \varphi, \psi) =
\text{length}_R(H^0(M, N, \varphi, \psi))
-
\text{length}_R(H^1(M, N, \varphi, \psi))
$$
In the case of a $(2, 1)$-periodic complex $(M, \varphi, \psi)$,
we denote this by $e_R(M, \varphi, \psi)$ and we will sometimes call this
the {\it (additive) Herbrand quotient}.
\end{definition}

\noindent
If the cohomology groups of $(M, \varphi, \psi)$
are finite abelian groups, then it is customary to call the
{\it (multiplicative) Herbrand quotient}
$$
q(M, \varphi, \psi) =
\frac{\# H^0(M, \varphi, \psi)}{\# H^1(M, \varphi, \psi)}
$$
In words: the multiplicative Herbrand quotient is the number of elements of
$H^0$ divided by the number of elements of $H^1$. If $R$ is local and if
the residue field of $R$ is finite with $q$ elements, then we see that
$$
q(M, \varphi, \psi) = q^{e_R(M, \varphi, \psi)}
$$

\medskip\noindent
An example of a $(2, 1)$-periodic complex over a ring $R$ is any triple of
the form $(M, 0, \psi)$ where $M$ is an $R$-module and $\psi$ is an
$R$-linear map. If the kernel and cokernel of $\psi$ have finite length,
then we obtain
\begin{equation}
\label{equation-multiplicity-coker-ker}
e_R(M, 0, \psi) = \text{length}_R(\Coker(\psi)) - \text{length}_R(\Ker(\psi))
\end{equation}
We state and prove the obligatory lemmas on these notations.

\begin{lemma}
\label{lemma-additivity-periodic-length}
Let $R$ be a ring. Suppose that we have a short exact sequence of
$2$-periodic complexes
$$
0 \to (M_1, N_1, \varphi_1, \psi_1)
\to (M_2, N_2, \varphi_2, \psi_2)
\to (M_3, N_3, \varphi_3, \psi_3)
\to 0
$$
If two out of three have cohomology modules of finite length so does
the third and we have
$$
e_R(M_2, N_2, \varphi_2, \psi_2) =
e_R(M_1, N_1, \varphi_1, \psi_1) +
e_R(M_3, N_3, \varphi_3, \psi_3).
$$
\end{lemma}

\begin{proof}
We abbreviate $A = (M_1, N_1, \varphi_1, \psi_1)$,
$B = (M_2, N_2, \varphi_2, \psi_2)$ and $C = (M_3, N_3, \varphi_3, \psi_3)$.
We have a long exact cohomology sequence
$$
\ldots
\to H^1(C)
\to H^0(A)
\to H^0(B)
\to H^0(C)
\to H^1(A)
\to H^1(B)
\to H^1(C)
\to \ldots
$$
This gives a finite exact sequence
$$
0 \to I
\to H^0(A)
\to H^0(B)
\to H^0(C)
\to H^1(A)
\to H^1(B)
\to K \to 0
$$
with $0 \to K \to H^1(C) \to I \to 0$ a filtration. By additivity of
the length function (Algebra, Lemma \ref{algebra-lemma-length-additive})
we see the result.
\end{proof}

\begin{lemma}
\label{lemma-finite-periodic-length}
Let $R$ be a ring. If $(M, N, \varphi, \psi)$ is a $2$-periodic complex
such that $M$, $N$ have finite length, then
$e_R(M, N, \varphi, \psi) = \text{length}_R(M) - \text{length}_R(N)$.
In particular, if $(M, \varphi, \psi)$ is a $(2, 1)$-periodic complex
such that $M$ has finite length, then
$e_R(M, \varphi, \psi) = 0$.
\end{lemma}

\begin{proof}
Observe that on the category of $2$-periodic complexes with $M$, $N$
of finite length the quantity ``$\text{length}_R(M) - \text{length}_R(N)$''
is additive in short exact sequences (precise statement left to the reader).
Consider the short exact sequence
$$
0 \to (M, \Im(\varphi), \varphi, 0) \to
(M, N, \varphi, \psi) \to (0, N/\Im(\varphi), 0, 0) \to 0
$$
The initial remark combined with the additivity of
Lemma \ref{lemma-additivity-periodic-length}
reduces us to the cases (a) $M = 0$ and (b) $\varphi$ is surjective.
We leave those cases to the reader.
\end{proof}

\begin{lemma}
\label{lemma-compare-periodic-lengths}
Let $R$ be a ring. Let $f : (M, \varphi, \psi) \to (M', \varphi', \psi')$
be a map of $(2, 1)$-periodic complexes whose cohomology modules
have finite length. If $\Ker(f)$ and $\Coker(f)$ have finite length,
then $e_R(M, \varphi, \psi) = e_R(M', \varphi', \psi')$.
\end{lemma}

\begin{proof}
Apply the additivity of Lemma \ref{lemma-additivity-periodic-length}
and observe that $(\Ker(f), \varphi, \psi)$ and
$(\Coker(f), \varphi', \psi')$ have vanishing multiplicity by
Lemma \ref{lemma-finite-periodic-length}.
\end{proof}




